% Part: first-order-logic
% Chapter: natural-deduction
% Section: proving-things

\documentclass[../../../include/open-logic-section]{subfiles}

\begin{document}

\iftag{FOL}
      {\olfileid{fol}{ntd}{pro}}
      {\olfileid{pl}{ntd}{pro}}

\olsection{\usetoken{P}{derivation}નાં ઉદાહરણો}

\begin{ex}
ચાલો, !!{sentence} $(!A \land !B) \lif !A$ની !!a{derivation} આપીએ.

!!{derivation}ના તળિયે ઇચ્છિત ફલિતવાક્ય લખીને શરૂ કરીએ.
\begin{prooftree}
\AxiomC{}
\UnaryInfC{$(!A\land !B) \lif !A$}
\end{prooftree}

હવે આ સ્વરૂપનું !!a{sentence} કયા પ્રકારના અનુમાનથી મળી શકે તે શોધવાનું
છે. ફલિતવાક્યનો !!{main operator} $\lif$ છે, તેથી \Intro{\lif} નિયમ
વાપરીને ફલિતવાક્ય સુધી પહોંચવાનો પ્રયત્ન કરીએ. આગળ વધતાં સંબંધિત
ધારણાઓ અને અનુમાન-નિયમોનાં ચિહ્નો લખતા જવું સારું, જેથી સાબિતીના અંતે
બધી ધારણાઓ !!{discharged} થઈ છે કે નહીં તે સહેલાઈથી જોઈ શકાય.
\begin{prooftree}
\AxiomC{$\Discharge{!A \land !B}{1}$}
\DeduceC{$!A$}
\DischargeRule{\Intro{\lif}}{1}
\UnaryInfC{$(!A\land !B) \lif !A$}
\end{prooftree}

હવે ધારણા $!A \land !B$થી $!A$ સુધીનાં પગલાં ભરવાનાં છે. અહીં માત્ર
એક સંયોજક, $\land$, સંભાળવાનો હોવાથી $\land$નો નિવારણ-નિયમ વાપરવો
જોઈએ. આમ નીચેની સાબિતી મળે છે:
\begin{prooftree}
\AxiomC{$\Discharge{!A \land !B}{1}$}
\RightLabel{\Elim{\land}}
\UnaryInfC{$!A$}
\DischargeRule{\Intro{\lif}}{1}
\UnaryInfC{$(!A\land !B) \lif !A$}
\end{prooftree}
હવે આપણી પાસે $(!A \land !B) \lif !A$ની યોગ્ય !!{derivation} છે.
\end{ex}

\begin{ex}
હવે $(\lnot !A \lor !B) \lif (!A \lif !B)$ની !!a{derivation} આપીએ.

!!{derivation}ના તળિયે ઇચ્છિત ફલિતવાક્ય લખીને શરૂ કરીએ.
\begin{prooftree}
\AxiomC{}
\UnaryInfC{$(\lnot !A \lor !B) \lif (!A \lif !B)$}
\end{prooftree}
આ ફલિતવાક્ય આપી શકે એવો તાર્કિક નિયમ શોધવા આપણે તેમાંના તાર્કિક
સંયોજકો જોઈએ છીએ: $\lnot$, $\lor$ અને $\lif$. અત્યારે આપણને $\lif$ના
પ્રથમ આવર્તનની જ ચિંતા છે, કારણ કે તે ફલિતવાક્યમાંના !!{sentence}નો
!!{main operator} છે; જ્યારે $\lnot$, $\lor$ અને $\lif$નું બીજું
આવર્તન બીજા સંયોજકના વ્યાપમાં છે, તેથી તેમને પછી સંભાળીશું. એટલે આપણે
\Intro{\lif} નિયમથી શરૂ કરીએ છીએ. તેનો યોગ્ય ઉપયોગ આવો હોવો જોઈએ:
\sourcecorrection{OLND-002}{સ્થિર અંગ્રેજી મૂળ આ સ્વાભાવિક-નિગમન
નિષ્પત્તિના અંતિમ વાક્યને ભૂલથી ``અંતિમ સિક્વન્ટમાંનું વાક્ય'' કહે છે;
અહીં આશય મુજબ તેને ફલિતવાક્યમાંનું વાક્ય કહ્યું છે.}
\begin{prooftree}
\AxiomC{$\Discharge{\lnot !A \lor !B}{1}$}
\DeduceC{$!A \lif !B$}
\DischargeRule{\Intro{\lif}}{1}
\UnaryInfC{$(\lnot !A \lor !B) \lif (!A \lif !B)$}
\end{prooftree}
હવે આગળ વધવાની બે શક્યતાઓ છે. કાં તો તળિયેથી ઉપર કામ ચાલુ રાખીને
\Intro{\lif} નિયમનો બીજો ઉપયોગ શોધી શકીએ, અથવા ઉપરથી નીચે કામ કરીને
\Elim{\lor} નિયમ લાગુ કરી શકીએ. ચાલો, બીજો રસ્તો લઈએ. ધારણા
$\lnot !A \lor !B$ને \Elim{\lor}નું સૌથી ડાબું આધારવાક્ય બનાવીએ.
\Elim{\lor}ના માન્ય ઉપયોગમાં બીજાં બે આધારવાક્યો ફલિતવાક્ય
$!A \lif !B$ સાથે તાદાત્મ્ય ધરાવતાં હોવાં જોઈએ, પરંતુ દરેકને અનુક્રમે
$\lnot !A \lor !B$ના બે વિયોજિત ઘટકોમાંથી એક એવી નવી ધારણામાંથી
નિષ્પન્ન કરી શકાય. તેથી આપણી !!{derivation} આવી દેખાશે:
\begin{prooftree}
\AxiomC{$\Discharge{\lnot !A \lor !B}{1}$}
\AxiomC{$\Discharge{\lnot !A}{2}$}
\DeduceC{$!A \lif !B$}
\AxiomC{$\Discharge{!B}{2}$}
\DeduceC{$!A \lif !B$}
\DischargeRule{\Elim{\lor}}{2}
\TrinaryInfC{$!A \lif !B$}
\DischargeRule{\Intro{\lif}}{1}
\UnaryInfC{$(\lnot !A \lor !B) \lif (!A \lif !B)$}
\end{prooftree}

જમણી બાજુની બંને શાખામાં આપણે $!A \lif !B$ !!{derive} કરવા માગીએ
છીએ; તે \Intro{\lif} વાપરીને કરવું ઉત્તમ છે.
\begin{prooftree}
\AxiomC{$\Discharge{\lnot !A \lor !B}{1}$}
\AxiomC{$\Discharge{\lnot !A}{2}, \Discharge{!A}{3}$}
\DeduceC{$!B$}
\DischargeRule{\Intro{\lif}}{3}
\UnaryInfC{$!A \lif !B$}
\AxiomC{$\Discharge{!B}{2}, \Discharge{!A}{4}$}
\DeduceC{$!B$}
\DischargeRule{\Intro{\lif}}{4}
\UnaryInfC{$!A \lif !B$}
\DischargeRule{\Elim{\lor}}{2}
\TrinaryInfC{$!A \lif !B$}
\DischargeRule{\Intro{\lif}}{1}
\UnaryInfC{$(\lnot !A \lor !B) \lif (!A \lif !B)$}
\end{prooftree}

!!{derivation}ના બાકી રહેલા બે ભાગોમાં વચ્ચે $\lnot !A$ અને $!A$માંથી,
અને ડાબે $!A$ અને $!B$માંથી, $!B$ની !!{derivation}s જોઈએ. પહેલા
પહેલો ભાગ લઈએ. $\lnot !A$ અને $!A$ એ \Elim{\lnot}નાં બે આધારવાક્યો છે:
\begin{prooftree}
\AxiomC{$\Discharge{\lnot !A}{2}$}
\AxiomC{$\Discharge{!A}{3}$}
\RightLabel{\Elim{\lnot}}
\BinaryInfC{$\lfalse$}
\DeduceC{$!B$}
\end{prooftree}
\FalseInt વાપરીને $!B$ ફલિતવાક્ય તરીકે મેળવી શાખા પૂરી કરી શકીએ.
\begin{prooftree}
\AxiomC{$\Discharge{\lnot !A \lor !B}{1}$}
\AxiomC{$\Discharge{\lnot !A}{2}$}
\AxiomC{$\Discharge{!A}{3}$}
\RightLabel{\Elim{\lnot}}
\BinaryInfC{$\lfalse$}
\RightLabel{\FalseInt}
\UnaryInfC{$!B$}
\DischargeRule{\Intro{\lif}}{3}
\UnaryInfC{$!A \lif !B$}
\AxiomC{$\Discharge{!B}{2}, \Discharge{!A}{4}$}
\DeduceC{$!B$}
\DischargeRule{\Intro{\lif}}{4}
\UnaryInfC{$!A \lif !B$}
\DischargeRule{\Elim{\lor}}{2}
\TrinaryInfC{$!A \lif !B$}
\DischargeRule{\Intro{\lif}}{1}
\UnaryInfC{$(\lnot !A \lor !B) \lif (!A \lif !B)$}
\end{prooftree}
\sourcecorrection{OLND-003}{આ નિષ્પત્તિમાં વ્યાઘાત આપતા દ્વિ-આધારી
પગલાને સ્થિર અંગ્રેજી મૂળ ભૂલથી $\Intro{\lfalse}$ કહે છે; ઉપરના
વર્ણન અને નિયમ મુજબ અહીં $\Elim{\lnot}$ મૂક્યું છે.}

હવે સૌથી જમણી શાખા જોઈએ. અહીં યાદ રાખવું જરૂરી છે કે !!{derivation}ની
વ્યાખ્યા ધારણાઓ \emph{નિવૃત્ત કરવાની છૂટ આપે છે}, પણ તેમ કરવું
\emph{ફરજિયાત કરતી નથી}. બીજા શબ્દોમાં, ધારણાઓ $!A$ અને $!B$માંથી
એકનો ઉપયોગ કર્યા વિના બીજામાંથી $!B$ નિષ્પન્ન કરી શકીએ તો તે યોગ્ય છે.
અને~$!B$માંથી $!B$ !!{derive} કરવું તુચ્છ છે: એકલું $!B$ એવી
!!a{derivation} છે અને કોઈ અનુમાનની જરૂર નથી. તેથી ધારણા~$!A$ને સીધી
કાઢી નાખી શકીએ છીએ.
\begin{prooftree}
\AxiomC{$\Discharge{\lnot !A \lor !B}{1}$}
\AxiomC{$\Discharge{\lnot !A}{2}$}
\AxiomC{$\Discharge{!A}{3}$}
\RightLabel{\Elim{\lnot}}
\BinaryInfC{$\lfalse$}
\RightLabel{\FalseInt}
\UnaryInfC{$!B$}
\DischargeRule{\Intro{\lif}}{3}
\UnaryInfC{$!A \lif !B$}
\AxiomC{$\Discharge{!B}{2}$}
\RightLabel{\Intro{\lif}}
\UnaryInfC{$!A \lif !B$}
\DischargeRule{\Elim{\lor}}{2}
\TrinaryInfC{$!A \lif !B$}
\DischargeRule{\Intro{\lif}}{1}
\UnaryInfC{$(\lnot !A \lor !B) \lif (!A \lif !B)$}
\end{prooftree}
નોંધો કે પૂરી થયેલી !!{derivation}માં સૌથી જમણું \Intro{\lif} અનુમાન
વાસ્તવમાં કોઈ ધારણાને નિવૃત્ત કરતું નથી.
\end{ex}

\begin{ex}
અત્યાર સુધી આપણને \FalseCl{} નિયમની જરૂર પડી નથી. તે વિશિષ્ટ છે,
કારણ કે તે નિયમના ફલિતવાક્યનું ઉપ-!!{formula} ન હોય એવી ધારણાને
નિવૃત્ત કરવાની છૂટ આપે છે. તેનો \FalseInt{} નિયમ સાથે નજીકનો સંબંધ
છે. હકીકતમાં, \FalseInt{} નિયમ એ \FalseCl{} નિયમનો વિશેષ કિસ્સો
છે---``અંતઃપ્રજ્ઞાવાદી તર્કશાસ્ત્ર'' નામનું એક તર્કશાસ્ત્ર છે જેમાં
માત્ર \FalseInt{} માન્ય છે. બીજું કંઈ કામ ન કરે ત્યારે \FalseCl{}
છેલ્લો ઉપાય છે. દાખલા તરીકે, માનો કે આપણે $!A \lor \lnot !A$
!!{derive} કરવા માગીએ છીએ. આપણી સામાન્ય વ્યૂહરચના $\Intro{\lor}$
વાપરીને $!A \lor \lnot !A$ !!{derive} કરવાનો પ્રયત્ન કરવાની હશે.
પરંતુ એ માટે કોઈ ધારણા વિના $!A$ અથવા $\lnot !A$માંથી એકને !!{derive}
કરવું પડે, જે શક્ય નથી. હવે \FalseCl{} મદદે આવે છે!
\begin{prooftree}
  \AxiomC{$\Discharge{\lnot(!A \lor \lnot !A)}{1}$}
  \DeduceC{$\lfalse$}
  \DischargeRule{\FalseCl}{1}
  \UnaryInfC{$!A \lor \lnot !A$}
\end{prooftree}
હવે આપણે $\lnot(!A \lor \lnot !A)$માંથી $\lfalse$ની !!a{derivation}
શોધીએ છીએ. $\lfalse$ એ $\Elim{\lnot}$નું ફલિતવાક્ય હોવાથી તેનો પ્રયત્ન કરીએ:
\begin{prooftree}
  \AxiomC{$\Discharge{\lnot(!A \lor \lnot !A)}{1}$}
  \DeduceC{$\lnot !A$}
  \AxiomC{$\Discharge{\lnot(!A \lor \lnot !A)}{1}$}
  \DeduceC{$!A$}
  \RightLabel{\Elim{\lnot}}
  \BinaryInfC{$\lfalse$}
  \DischargeRule{\FalseCl}{1}
  \UnaryInfC{$!A \lor \lnot !A$}
\end{prooftree}
~$\lnot !A$ની !!a{derivation} શોધવાની આપણી વ્યૂહરચના
~$\Intro{\lnot}$નો ઉપયોગ કહે છે:
\begin{prooftree}
  \AxiomC{$\Discharge{\lnot(!A \lor \lnot !A)}{1}, \Discharge{!A}{2}$}
  \DeduceC{$\lfalse$}
  \DischargeRule{\Intro{\lnot}}{2}
  \UnaryInfC{$\lnot !A$}
  \AxiomC{$\Discharge{\lnot(!A \lor \lnot !A)}{1}$}
  \DeduceC{$!A$}
  \RightLabel{\Elim{\lnot}}
  \BinaryInfC{$\lfalse$}
  \DischargeRule{\FalseCl}{1}
  \UnaryInfC{$!A \lor \lnot !A$}
\end{prooftree}
અહીં ધારણા $\lnot(!A \lor \lnot !A)$ અને આપણી નવી ધારણા $!A$માંથી
~$\Intro{\lor}$ વડે મળતા $!A \lor \lnot !A$ પર $\Elim{\lnot}$ લગાડીને
$\lfalse$ સહેલાઈથી મેળવી શકીએ:
\begin{prooftree}
  \AxiomC{$\Discharge{\lnot(!A \lor \lnot !A)}{1}$}
  \AxiomC{$\Discharge{!A}{2}$}
  \RightLabel{\Intro{\lor}}
  \UnaryInfC{$!A \lor \lnot !A$}
  \RightLabel{\Elim{\lnot}}
  \BinaryInfC{$\lfalse$}
  \DischargeRule{\Intro{\lnot}}{2}
  \UnaryInfC{$\lnot !A$}
  \AxiomC{$\Discharge{\lnot(!A \lor \lnot !A)}{1}$}
  \DeduceC{$!A$}
  \RightLabel{\Elim{\lnot}}
  \BinaryInfC{$\lfalse$}
  \DischargeRule{\FalseCl}{1}
  \UnaryInfC{$!A \lor \lnot !A$}
\end{prooftree}
જમણી બાજુએ એ જ વ્યૂહરચના વાપરીએ છીએ, ફક્ત~$!A$ને~\FalseCl{}થી મેળવીએ છીએ:
\begin{prooftree}
  \AxiomC{$\Discharge{\lnot(!A \lor \lnot !A)}{1}$}
  \AxiomC{$\Discharge{!A}{2}$}
  \RightLabel{\Intro{\lor}}
  \UnaryInfC{$!A \lor \lnot !A$}
  \RightLabel{\Elim{\lnot}}
  \BinaryInfC{$\lfalse$}
  \DischargeRule{\Intro{\lnot}}{2}
  \UnaryInfC{$\lnot !A$}
  \AxiomC{$\Discharge{\lnot(!A \lor \lnot !A)}{1}$}
  \AxiomC{$\Discharge{\lnot !A}{3}$}
  \RightLabel{\Intro{\lor}}
  \UnaryInfC{$!A \lor \lnot !A$}
  \RightLabel{\Elim{\lnot}}
  \BinaryInfC{$\lfalse$}
  \DischargeRule{\FalseCl}{3}
  \UnaryInfC{$!A$}
  \RightLabel{\Elim{\lnot}}
  \BinaryInfC{$\lfalse$}
  \DischargeRule{\FalseCl}{1}
  \UnaryInfC{$!A \lor \lnot !A$}
\end{prooftree}
\end{ex}

\begin{prob}
નીચેનું દર્શાવતી !!{derivation}s આપો:
\begin{enumerate}
\item $!A \land (!B \land !C) \Proves (!A \land !B) \land !C$.
\item $!A \lor (!B \lor !C) \Proves (!A \lor !B) \lor !C$.
\item $!A \lif (!B \lif !C) \Proves !B \lif (!A \lif !C)$.
\item $!A \Proves \lnot\lnot !A$.
\end{enumerate}
\end{prob}

\begin{prob}
નીચેનું દર્શાવતી !!{derivation}s આપો:
\begin{enumerate}
\item $(!A \lor !B) \lif !C \Proves !A \lif !C$.
\item $(!A \lif !C) \land (!B \lif !C) \Proves (!A \lor !B) \lif !C$.
\item $\Proves \lnot(!A \land \lnot !A)$.
\item $!B \lif !A \Proves \lnot !A \lif \lnot !B$.
\item $\Proves (!A \lif \lnot !A) \lif \lnot !A$.
\item $\Proves \lnot(!A \lif !B) \lif \lnot !B$.
\item $!A \lif !C \Proves \lnot (!A \land \lnot !C)$.
\item $!A \land \lnot !C \Proves \lnot (!A \lif !C)$.
\item $!A \lor !B, \lnot !B \Proves !A$.
\item $\lnot !A \lor \lnot !B \Proves \lnot(!A \land !B)$.
\item $\Proves (\lnot !A \land \lnot !B) \lif\lnot(!A \lor !B)$.
\item $\Proves \lnot(!A \lor !B) \lif (\lnot !A \land \lnot !B)$.
\end{enumerate}
\end{prob}

\begin{prob}
નીચેનું દર્શાવતી !!{derivation}s આપો:
\begin{enumerate}
\item $\lnot(!A \lif !B) \Proves !A$.
\item $\lnot(!A \land !B) \Proves \lnot !A \lor \lnot !B$.
\item $!A \lif !B \Proves \lnot !A \lor !B$.
\item $\Proves \lnot \lnot !A \lif !A$.
\item $!A \lif !B, \lnot !A \lif !B \Proves !B$.
\item $(!A \land !B) \lif !C \Proves (!A \lif !C) \lor (!B \lif !C)$.
\item $(!A \lif !B) \lif !A \Proves !A$.
\item $\Proves (!A \lif !B) \lor (!B \lif !C)$.
\end{enumerate}
(આ બધામાં $\FalseCl$~નિયમ જરૂરી છે.)
\end{prob}

\end{document}
